% Generated draft: requires mathematical and source review.
\documentclass[11pt]{article}
\usepackage[margin=25mm]{geometry}
\usepackage{fontspec}
\setmainfont{Latin Modern Roman}
\usepackage{amsmath,amssymb,mathtools,unicode-math}
\setmathfont{Latin Modern Math}
\usepackage{enumitem,needspace,xurl,hyperref}
\setlength{\parindent}{0pt}
\setlength{\parskip}{5pt}
\setlength{\emergencystretch}{3em}
\newcommand{\N}{\mathbb{N}}
\newcommand{\Z}{\mathbb{Z}}
\newcommand{\Q}{\mathbb{Q}}
\newcommand{\R}{\mathbb{R}}
\newcommand{\C}{\mathbb{C}}
\newcommand{\F}{\mathbb{F}}
\newcommand{\A}{\mathbb{A}}
\newcommand{\PP}{\mathbb{P}}
\newcommand{\E}{\mathbb{E}}
\newcommand{\eps}{\varepsilon}
\newcommand{\dd}{\,\mathrm{d}}
\begin{document}
\section*{UnsolvedMath ID: 30001844}
\textbf{Source problem code:} catalogue entry
\subsection{Definitions and mathematical statement}
Let \(K\) be a field. Under the usual projective interpretation, a smooth cubic surface over \(K\) is a smooth degree-three hypersurface \(S\subseteq\PP^3_K\). Write \(S(K)\) for its set of \(K\)-rational points.

The supplied record does not define the permitted secant and tangent constructions completely. One precise finite-intersection convention is the following. Given \(P,Q\in S(K)\), possibly equal, a construction produces \(R\in S(K)\) if there is a line \(\ell\subseteq\PP^3_K\), defined over \(K\) and not contained in \(S\), whose intersection divisor satisfies
\[
\ell\cdot S=P+Q+R.
\]
Intersection multiplicities are included. For \(P\ne Q\), this is the secant construction. For \(P=Q\), the line has intersection multiplicity at least two at \(P\), giving the tangent construction.

For \(B\subseteq S(K)\), let \(\operatorname{cl}_{\mathrm{st}}(B)\) be the smallest subset of \(S(K)\) containing \(B\) and closed under these constructions. Thus its points are obtainable by finitely many successive constructions.

\textbf{Question, under this convention.} Is it true that, for every field \(K\) and every smooth cubic surface \(S\subseteq\PP^3_K\), there exists a finite subset \(B\subseteq S(K)\) such that
\[
\operatorname{cl}_{\mathrm{st}}(B)=S(K)?
\]
The quantification over all fields preserves the supplied wording; no restriction to number fields, finitely generated fields, infinite fields, or particular characteristics is supplied. If \(S(K)=\varnothing\), the choice \(B=\varnothing\) satisfies the displayed equality.

\textbf{Unresolved qualifications.} The projective hypersurface interpretation and the finite-intersection convention above make the question explicit, but the supplied record does not confirm these conventions. In particular, it does not specify whether constructions using lines contained in \(S\) may add points on those lines. Such lines have no residual degree-three intersection divisor. The intended construction rule and intended field hypotheses require source review before this draft can be treated as a definitive formulation.
\subsection{Short English statement}
The supplied question asks whether finitely many rational points on every smooth cubic surface over a field generate all its rational points by repeated secant and tangent constructions. The draft specifies a residual-intersection interpretation, while flagging the unspecified treatment of lines contained in the surface and the unrestricted field domain. The supplied status is \emph{partially solved}; this is a curation label, not an independently verified conclusion.
\subsection{Sources}
\begin{itemize}
\item Supplied local UnsolvedMath registry record, permanent ID \(30001844\), title \emph{Finite Secant--Tangent Generation on Cubic Surfaces}; Oberwolfach Reports Open Problems collection; exact source problem code \texttt{catalogue entry}. The record gives proposed year \(2011\) and names no proposer.
\item Supplied report citation: \emph{Explicit Methods in Number Theory (2012)}, Oberwolfach Reports, pp.\ 1971--2044, DOI \texttt{10.4171/owr/2011/35}. \url{https://doi.org/10.4171/owr/2011/35}. The background describes Oberwolfach Reports as documenting mathematical workshops and their open-problem discussions. The citation year \(2012\), proposed year \(2011\), and DOI containing \(2011\) are retained as supplied.
\item Supplied literature reference: ``Siksek on Mordell--Weil generators.'' \url{https://arxiv.org/abs/1012.1838}.
\item Supplied literature reference: ``Manin's cubic-surface problem.'' \url{https://arxiv.org/abs/alg-geom/9407009}.
\item Literature assessment dated exactly \texttt{2026-08-21}: the supplied curation states that Manin's Mordell--Weil problem for cubic surfaces remains open in general and that Siksek proves a one-generator result for surfaces containing two skew rational lines. This is supplied dated curation derived from the OWR record, not an independently verified literature assessment. The record does not supply the precise field hypotheses or other hypotheses of the cited one-generator result.
\item The supplied statement status is \texttt{corrected\_verified}. Its accompanying verification text reports comparison with the official Oberwolfach report and normalization of notation, context, and question form without changing the source-level claim. This is a supplied verification assertion; no source consultation or independent mathematical verification is claimed here.
\end{itemize}
\end{document}
